\documentclass[main]{subfiles}
\usepackage[sols]{workbook}
\numbering{ws2-6}

\begin{document}

\ifSubfilesClassLoaded{%
  \chaptermark{\chaptertwo}
}{}

\section{Geometric Series} \label{ws2-6}
\begin{framed}
\centerline{\textbf{Lesson Goals}}

You should be able to:
\begin{itemize}
\item Determine whether a given series is geometric, and identify the first term and the common ratio for a geometric series
\item Find what value a convergent geometric series converges to 
\item Determine if a given geometric series converges or diverges
\end{itemize}
\end{framed}

 
 \begin{enumerate}
     

    \item If you suffer from allergies, your doctor may
    suggest that you take Claritin once a day.  Each Claritin tablet
    contains 10 mg of loratadine (the active ingredient).  Because of the
    way the human body metabolizes loratadine, no matter how much is in your
    body at a given time, only $\frac{1}{8}$ of that amount remains 24 hours
    later.\footnote{This estimate comes from the fact that the average
    half-life of loratadine is known to be 8 hours.}
    \begin{enumerate}
      \item
      \begin{problem}[\vfill]
      If you take one Claritin tablet each morning for 2 successive days,
      how much loratadine is in your body right after you take the 2nd
      tablet? (\textbf{Don't} simplify your answer! Just write an arithmetic expression that gives the amount.)
      \end{problem}

      \begin{solution}
      Immediately after taking the first tablet, you have 10 mg of
      loratadine in your body.  The following morning, only $\frac{1}{8}$
      of that is left, so you have $10 \left( \frac{1}{8} \right)$ mg in
      your body.  You then take another pill containing 10 mg, so you have
      a total of $\boxed{10 + 10 \left( \frac{1}{8} \right)}$ mg of
      loratadine in your body.
      \end{solution}      

      \item
      \begin{problem}[\vfill\vfill]
      If you take one Claritin tablet every morning for a week, how much
      loratadine is in your body right after you take the 3rd tablet?  7th
      tablet?  (\textbf{Don't} try to simplify your computations! Just write out an
      arithmetic expression that gives the amount. Try to find a pattern.)
      \end{problem}

      \begin{solution}
      In (a), we established that you have $10 + 10 \left( \frac{1}{8}
      \right)$ mg of loratadine in your body immediately after taking the
      2nd tablet.

      The following morning, 1/8 of that remains, or $\left[10 + 10 \left(
      \frac{1}{8} \right) \right] \frac{1}{8} = 10 \left( \frac{1}{8}
      \right) + 10 \left( \frac{1}{8} \right)^2$.  You then take another
      pill containing 10 mg, so you have a total of $10 + 10 \left(
      \frac{1}{8} \right) + 10 \left( \frac{1}{8} \right)^2$ mg of
      loratadine in your body after taking the 3rd pill.

      Continuing this reasoning, you will have $10 + 10 \left( \frac{1}{8}
      \right) + 10 \left( \frac{1}{8} \right)^2+ 10 \left( \frac{1}{8}
      \right)^3 + 10 \left( \frac{1}{8} \right)^4+ 10 \left( \frac{1}{8}
      \right)^5 + 10 \left( \frac{1}{8} \right)^6$ mg in your body after
      the 7th pill.
      \end{solution}

      \item
      \begin{problem}[\vfill\vfill\vfill]
      If you take Claritin for years and years, will the amount of
      loratadine in your body level off?  Or will your bloodstream be pure
      loratadine?
      \end{problem}

      \begin{solution}
      Right after you take the $n$-th pill, you will have $10 + 10 \left(
      \frac{1}{8} \right) + 10 \left( \frac{1}{8} \right)^2 + 10 \left(
      \frac{1}{8} \right)^3 + \cdots + 10 \left( \frac{1}{8}
      \right)^{n-1}$ mg of loratadine in your body.  Let's call this
      amount $s_n$.  We are wondering what happens to $s_n$ as $n$ gets
      very large.

      We will use a clever algebraic manipulation to get a closed form expression for $s_n$.  We
      know
      \begin{equation}
      s_n = 10 + 10 \left( \frac{1}{8} \right) + 10 \left( \frac{1}{8}
      \right)^2 + \cdots + 10 \left( \frac{1}{8} \right)^{n-2} + 
      10 \left( \frac{1}{8} \right)^{n-1}
      \end{equation}
      Notice that, if we multiply both sides by $\frac{1}{8}$ (multiplying
      term-by-term on the right side), we get something similar looking on
      the right side:
      \begin{equation}
      \frac{1}{8} s_n = 10 \left(
      \frac{1}{8} \right) + 10 \left( \frac{1}{8} \right)^2 + 10
      \left( \frac{1}{8} \right)^3 + \cdots + 10 \left( \frac{1}{8}
      \right)^{n-1} + 10 \left( \frac{1}{8} \right)^n
      \end{equation}
      If we subtract (2) from (1), most of the terms on
      the right side cancel, and we are left with $\displaystyle
      \frac{7}{8} s_n = 10 - 10 \left( \frac{1}{8} \right)^n$.  Dividing
      both sides by $\displaystyle \frac{7}{8}$, $\displaystyle s_n =
      \frac{10 - 10(1/8)^n}{7/8}$.

      Using this expression for $s_n$, it is easy to see what happens as
      $n$ gets bigger and bigger: $\left( \frac{1}{8} \right)^n$ gets
      closer and closer to 0, so $\frac{10 - 10(1 / 8)^n}{7 / 8}$ gets
      closer and closer to $\frac{10}{7 / 8} = \frac{80}{7}$.  Thus, over
      time, the amount of loratadine in your body gets closer and closer
      to $\frac{80}{7}$ mg.

      Notice that what we've shown mathematically is that the sum of the
      infinite series $10 + 10 \left( \frac{1}{8} \right) + 10 \left(
      \frac{1}{8} \right)^2 + \cdots$ is $\frac{80}{7}$.  After all,
      $s_n$ is just the $n$-th partial sum of this series, and we showed
      that $\displaystyle \lim_{n \to \infty} s_n = \frac{80}{7}$.
      \end{solution}

    \end{enumerate}
  \end{enumerate}
  \pspace{\newpage}
 \begin{problemonly}
  \begin{framed}
  \textbf{Definition.} A \uline{geometric series} is a series that can be
  written in the form $\displaystyle \sum_{k=0}^\infty ar^k = a + ar +
  ar^2 + ar^3 + \cdots$.  In particular, a geometric series has a \textit{common ratio} $r$ such that any term multiplied by $r$ gives the next term in the series.
  \end{framed}  
  \end{problemonly}

  \begin{enumerate}[resume]

    \item  Which of the
  following series are geometric?  For each geometric series, identify $a$
  and $r$. 
  
  \textbf{Hint}: Start by writing out the first few terms of each series in $\cdots$ notation!
  \begin{enumerate}
    \item
    \begin{problem}[\vfill]
    $\displaystyle \sum_{k=0}^\infty \frac{(-1)^k 4^{k+2}}{3^{2k}}$.
    \end{problem}

    \begin{solution}
This series is geometric. $a=16$ and $r=\frac{-4}{9}$.
    \end{solution}

    \item
    \begin{problem}[\vfill]
    $\displaystyle \sum_{k=1}^{\infty} \frac{2^k}{k^3}$.
    \end{problem}

    \begin{solution}
    If we write this series out, it is $2 + \frac{4}{8} + \frac{8}{27} +
    \frac{16}{64} + \cdots$.  We can already see that it's not geometric
    because the terms don't have a common ratio.  (To elaborate: if it
    was geometric, the first term would be $a$ and the second term would
    be $ar$; this means that $a$ would have to be 1, and $r$ would have
    to be $\frac{1}{8}$.  But then the third term isn't $\frac{1}{8}$ times the second term,
    so there is no geometric series pattern here.)
    \end{solution}

    \item
    \begin{problem}[\vfill]
    $\displaystyle \sum_{n=4}^{\infty} (-1)^n\left(\frac{2}{3}\right)^n$.
    \end{problem}

    \begin{solution}
   This series is geometric. $a=\left(\frac{2}{3}\right)^4=\frac{16}{81}$ and $r=\frac{-2}{3}$.
    \end{solution}
\end{enumerate}
  \pspace{\newpage}
  \item 
  Let $s_n$ be the $n$-th partial sum
  of the geometric series $a + ar + ar^2 + ar^3 + \cdots$.  That is, $s_n$
  is the sum of the first $n$ terms.  In this problem, we'll consider only
  geometric series with $a \neq 0$.
  \begin{enumerate}

    \item
    \begin{problem}[0.75in]
    Write a closed form expression for $s_n$ if $r = 1$.  (A ``closed
    form expression'' is a formula without $\cdots$ or $\sum$; if you had
    values for $a$ and $n$, you could plug them into your formula 
    and evaluate it directly.)
    \end{problem}

    \begin{solution}
    If $r = 1$, then $s_n = a + a + \cdots + a$ (with $n$ terms), so
    $s_n = na$. 
    \end{solution}
    
    \item
    \begin{problem}[2.5in]
    Write a closed form expression for $s_n$, assuming $r \neq 1$.
    \end{problem}

    \begin{solution}
    Since $s_n$ is the sum of the first $n$ terms of the series, $s_n =
    a + ar + ar^2 + \cdots + ar^{n-1}$.  We multiply this equation by
    $r$ and then subtract that from the original equation:
    \[ \begin{array}{rcrcrcrcccrcr}
    s_n & = & a & + & ar & + & ar^2 & + & \cdots & + & ar^{n-1}
    \\
    r s_n & = & & & ar & + & ar^2 & + & \cdots & + &
    ar^{n-1} & + & ar^n \\
    \hline 
    (1 - r) s_n & = & a & + & 0 & + & 0 & + & \cdots & + & 0 & - & a
    r^n 
    \end{array}\]
    Dividing both sides by $1 - r$, $s_n = \boxed{\frac{a - ar^n}{1 -
    r}}$.  (It's fine to divide by $1 - r$ since $r \neq 1$.)
    \end{solution}

    \item
    \begin{problem}[3in]
    For what values of $a$ and $r$ does the geometric series $a + ar +
    ar^2 + \cdots$ converge?  For what values does it diverge?  When it
    converges, what is its sum?

    \probonly{\textbf{This result is incredibly important; you should
    feel free to cite it from now on without proof.  When using it, be sure to
    always explain what $a$ (the first term) and $r$ (the common ratio) are for your geometric series.}}
    \end{problem}

    \begin{solution}
    By definition of convergence, we need to look at the limit of
    partial sums $\displaystyle \lim_{n \to \infty} s_n$.  If $r \neq
    1$, we have shown that this is equal to $\displaystyle \lim_{n \to
    \infty} \frac{a - ar^n}{1 - r} = \frac{a}{1-r}\cdot\lim_{n \to \infty} (1-r^n)$. (We can factor out $\frac{a}{1-r}$, since the limit only depends on $n$.)  When $|r| > 1$ or $r = -1$, this
    limit does not exist (if $|r| > 1$, then $r^n$ just grows bigger and
    bigger in magnitude; if $r = -1$, then $r^n$ just bounces between
    $1$ and $-1$ forever).  When $|r| < 1$, $r^n$ tends to 0, so
    $\displaystyle \lim_{n \to \infty} \frac{a - ar^n}{1 - r} =
    \frac{a}{1 - r}$.

    When $r = 1$, we have shown that $\displaystyle \lim_{n \to \infty}
    s_n = \lim_{n \to \infty} na$, which does not exist.

    So, to summarize,
   \[ a + ar + ar^2 + ar^3 + \cdots \left\{
    \begin{array}{l l}
    \textrm{converges to } \frac{a}{1 - r} & \textrm{ if } |r| < 1 \\
    \textrm{diverges} & \textrm{ if } |r| \geq 1
    \end{array} \right.\]
    \end{solution}
  \end{enumerate}
  \pspace{\newpage}
  \item 
  Remember that a geometric series is a series that can be expressed in the form $a + ar +
  ar^2 + \cdots = \displaystyle \sum_{k=0}^\infty ar^k$.  Which of the
  following series are geometric?  For each geometric series, identify $a$
  and $r$, and say whether the series converges or diverges; if it
  converges, find its sum. 
  
  \textbf{Hint}: Start by writing out the first few terms of each series in $\cdots$ notation!
  \begin{enumerate}
    \item
    \begin{problem}[\vfill]
    $\displaystyle \sum_{k=1}^\infty \frac{(-1)^k 2^{k+1}}{3^k}$.
    \end{problem}

    \begin{solution}
    We can write the given series in $\cdots$ notation as $- \frac{4}{3}
    + \frac{8}{9} - \frac{16}{27} + \cdots$.  If this is geometric, then
    $a = - \frac{4}{3}$ and $ar = \frac{8}{9}$, so $r = - \frac{2}{3}$.

    To see whether the series really is geometric, we need to make sure
    that the rest of the series follows the pattern that each term is
    $r$ times the previous term.

    If we write the $k$-th term as $a_k = \frac{(-1)^k 2^{k+1}}{3^k}$,
    then the next term is $a_{k+1} = \frac{(-1)^{k + 1} 2^{k+2}}{3^{k +
    1}}$, which is indeed equal to $-\frac{2}{3} a_k$.  So, the
    series is geometric with $a = -\frac{4}{3}$ and $r = -\frac{2}{3}$.
    Since $|r| < 1$, the series converges, and its sum is $\frac{a}{1 -
    r} = -\frac{4}{5}$.
    \end{solution}

    \item
    \begin{problem}[\vfill]
    $\displaystyle \sum_{k=1}^{\infty} \frac{1}{k^3}$.
    \end{problem}

    \begin{solution}
    If we write this series out, it is $1 + \frac{1}{8} + \frac{1}{27} +
    \frac{1}{64} + \cdots$.  We can already see that it's not geometric
    because the terms don't have a common ratio.  (To elaborate: if it
    was geometric, the first term would be $a$ and the second term would
    be $ar$; this means that $a$ would have to be 1, and $r$ would have
    to be $\frac{1}{8}$.  But then the third term isn't $\frac{1}{8}$ times the second term,
    so there is no geometric series pattern here.)
    \end{solution}

    \item
    \begin{problem}[\vfill]
    $\displaystyle \sum_{n=4}^{\infty} \frac{2}{3^{2n}}$.
    \end{problem}

    \begin{solution}
    $\displaystyle \sum_{n=4}^{\infty} \frac{2}{3^{2n}} = \frac{2}{3^8}
    + \frac{2}{3^{10}} + \frac{2}{3^{12}} + \cdots$, which looks like it
    could be geometric with $a = \frac{2}{3^8}$ and $r =
    \frac{1}{3^2}$.  Indeed, if we write $a_n = \frac{2}{3^{2n}}$, then
    the next term is $a_{n+1} = \frac{2}{3^{2 (n+ 1)}} = \frac{2}{3^{2n
    + 2}}$, which is equal to $\frac{1}{3^2} a_n$, so each term is
    $r$ times the previous.

    So, the series is geometric with $a = \frac{2}{3^8}$ and $r =
    \frac{1}{3^2}$.  Since $|r| < 1$, the series converges and has sum
    $\frac{a}{1 - r} = \frac{1}{2916}$.
    \end{solution}
\end{enumerate} 

\pspace{\newpage}
  \item 
  \begin{enumerate}
    \item
    \begin{problem}[\vfill]
    For what values of $x$ does the geometric series $\displaystyle
    \sum_{k=0}^\infty x^k = 1 + x + x^2 + \cdots$ converge?  For what
    values of $x$ does it diverge?  When it converges, what does it
    converge to?
    \end{problem}
    
    \begin{solution}
    This is a geometric series $a + ar + ar^2 + ar^3 + \cdots$ with $a =
    1$ and $r = x$ (remember that when we write a geometric series as $a
    + ar + ar^2 + ar^3 + \cdots$, this means that the first term is $a$
    and the common ratio is $r$).  Therefore, we know that it diverges
    when $|x| \geq 1$.  When $|x| < 1$, it converges to $\displaystyle \frac{a}{1 - r}
    = \frac{1}{1 - x}$.
    \end{solution}

    \item
    \begin{problem}[\vfill]
    How would you graph the function $\displaystyle 1+x+x^2+x^3+\cdots$?
    \end{problem}
    
    \begin{solution}
    We know from {{\#3}(a)} that 
 \[ \displaystyle 1+x+x^2+\cdots=\sum_{k=0}^\infty x^k= 
    \begin{cases}
   \frac{1}{1 - x} &\text{ when } |x| < 1 \\
    \text{diverges} &\text{ when } |x| \geq 1
    \end{cases}\]
   
    So, when $|x| < 1$, the graph of $\displaystyle \sum_{k=0}^\infty
    x^k$ should look exactly like the graph of $\frac{1}{1 - x}$.  When
    $|x| \geq 1$, $\displaystyle \sum_{k=0}^\infty x^k$ diverges, which
    basically means it's undefined.  So, $\displaystyle
    \sum_{k=0}^\infty x^k$ really has a domain of just $(-1,1)$, and on
    that domain, it should look exactly like $\frac{1}{1 - x}$.  So, we
    end up with this graph of $\displaystyle \sum_{k=0}^\infty x^k$: 
    \begin{center}
    \includegraphics[width=2in]{ws/images/geometric-power-series-1}
    \end{center}
    For reference, here is the graph of $\frac{1}{1 - x}$:
    \begin{center}
    \includegraphics[width=2in]{ws/images/geometric-power-series-1-base}
    \end{center}
    
    Most importantly, recall that $1+x+^2+\cdots$ is the Taylor series for the function $f(x)=\frac{1}{1-x}$ centered at 0. Therefore, this result also tells us that the Taylor series for $f(x)=\frac{1}{1-x}$ centered at 0 is actually \textit{equal} to $f(x)$ on the interval $(-1.1)$. This is the first example we have found where we can prove that a Taylor series of a function is not merely an approximation of the function, but actually equivalent to the function for some $x$-values!
    \end{solution}
   
  \end{enumerate}
  

 
  

      \item 
  \begin{problem}[\vfill]
Use a geometric series to show that $0.999...=1$. 

%\textbf{Hint:} How can you express $0.999...$ as a sum of infinitely many fractions? What is the first term $a$? What is the common ratio $r$?
  \end{problem}
  
  \begin{solution}
  The repeating decimal $0.999...$ can be expressed as:
  
  $$0.999...=0.9+0.09+0.009+\cdots=\frac{9}{10}+\frac{9}{100}+\frac{9}{1000}+\cdots$$
  
  This is just a geometric series whose first term is $a=\frac{9}{10}$ and whose common ratio is $r=\frac{1}{10}$. Since $|r|<1$, this geometric series converges, and thus its sum is given by $\frac{a}{1-r}$. Therefore,
  
  $$0.999...=\frac{9}{10}+\frac{9}{100}+\frac{9}{1000}+\cdots=\frac{9/10}{1-1/10}=\frac{9/10}{9/10}=1$$
  
  Many people struggle to believe this, thinking that 1 must be larger than $0.999...$. (More generally, it can be hard to believe that a convergent infinite sum is actually \textit{equal} to the number it converges to, instead of just \textit{really close to} the number to which it converges.) 
  
However, if you believe that 1 is larger than $0.999...$, then it must be larger by some finite positive amount. But what is that amount?
  
  \end{solution}
 
  


    
   

  \end{enumerate}

\end{document}